% ----------------------------------------------------------------------------
%  6. Methodology
% ----------------------------------------------------------------------------
\section{Methodology}
\label{sec:methodology}

\noindent
The engine has four pluggable pieces. Each can be swapped without
touching the others.

\subsection{Ordering --- ``what task next?''}

\begin{itemize}
  \item \textbf{FCFS.}\; First-come-first-served queue;
        ties broken by task ID.
  \item \textbf{Priority.}\; Min-heap on a priority key;
        smaller key = more urgent.
\end{itemize}

\subsection{Matching --- ``who does the task?''}

\begin{itemize}
  \item \textbf{Greedy.}\; Pick the role-matched worker who is
        free earliest. Can stack many tasks on one worker per
        tick.
  \item \textbf{FairGreedy($\alpha$).}\;
        Trade off ``free earliest'' vs.\ ``most shift remaining.''
        $\alpha\!=\!1$ is plain Greedy; $\alpha\!=\!0$ is pure
        load-balancing.
  \item \textbf{Hungarian.}\; Solve the optimal one-to-one match
        between this tick's ready tasks and free workers.
        Overflow tasks are deferred to the next tick.
\end{itemize}

\subsection{Priority key --- the tie-breaker family}

\begin{itemize}
  \item Each policy expresses urgency as a list of integers
        compared left-to-right.
  \item Smaller is more urgent; negation flips a single dimension.
  \item Default reproduces the legacy
        $\text{acuity}\!\times\!10 + \text{priority}$ formula.
  \item Last dimension is a seeded hash so equal-key ties resolve
        deterministically.
\end{itemize}

\begin{table}[H]
\centering
\small
\caption{The seven priority policies the system ships with.}
\label{tab:lex_policies}
\setlength{\tabcolsep}{6pt}
\begin{tabularx}{\textwidth}{l X}
\toprule
\textbf{Policy} & \textbf{Key shape (smaller better on every dim)} \\
\midrule
\code{acuity\_first}        & (acuity, base-priority, hash) \\
\code{priority\_first}      & (base-priority, acuity, hash) \\
\code{acuity\_short\_first} & (acuity, $+$duration, base-priority, hash) \\
\code{acuity\_long\_first}  & (acuity, $-$duration, base-priority, hash) \\
\code{acuity\_old\_first}   & (acuity, $-$age, base-priority, hash) \\
\code{acuity\_young\_first} & (acuity, $+$age, base-priority, hash) \\
\code{duration\_first}      & ($+$duration, acuity, base-priority, hash) \\
\bottomrule
\end{tabularx}
\end{table}

\subsection{Picking the best schedule from many}

\begin{itemize}
  \item Each candidate schedule has three ``smaller is better''
        scores:
        \begin{itemize}
          \item Urgent-care: average finish time of CRIT+HIGH
                patients.
          \item Makespan: when the schedule ends overall.
          \item Worst overtime: most overtime any single worker
                takes.
        \end{itemize}
  \item Drop dominated candidates (some other candidate is
        $\leq$ on every score and $<$ on at least one).
  \item Among the survivors, pick six spread out across the
        three axes so the human sees the trade-offs, not just one
        corner.
\end{itemize}

\subsection{Capacity feasibility}

\begin{itemize}
  \item If a task pushes a worker past their shift, the engine
        records the overtime but still schedules the task.
  \item Post-run: total overtime, per-worker overtime, and a
        boolean ``feasible at threshold $\theta$.''
  \item The engine never refuses to schedule on overtime grounds
        --- the human is the final judge.
\end{itemize}
